\documentclass[11pt,a4paper]{article}

% ---------------------------------------------------------------------------
% Hopper -- Local Setup Journey & Architecture Reference
% Compile with:  pdflatex hopper-setup-and-architecture.tex  (run twice for ToC)
% ---------------------------------------------------------------------------

\usepackage[utf8]{inputenc}
\usepackage[T1]{fontenc}
\usepackage[margin=1in]{geometry}
\usepackage{xcolor}
\usepackage{listings}
\usepackage{booktabs}
\usepackage{longtable}
\usepackage{enumitem}
\usepackage{parskip}
\usepackage{titlesec}
\usepackage[hidelinks]{hyperref}
\usepackage{fancyhdr}

% ---- Colours ----
\definecolor{codebg}{gray}{0.96}
\definecolor{codeframe}{gray}{0.80}
\definecolor{kw}{RGB}{0,90,160}
\definecolor{cmt}{RGB}{100,120,100}
\definecolor{str}{RGB}{160,60,40}
\definecolor{accent}{RGB}{20,80,140}

% ---- Listings ----
\lstdefinestyle{shell}{
  basicstyle=\ttfamily\footnotesize,
  backgroundcolor=\color{codebg},
  frame=single, rulecolor=\color{codeframe},
  breaklines=true, breakatwhitespace=false,
  columns=fullflexible, keepspaces=true,
  showstringspaces=false,
  keywordstyle=\color{kw}\bfseries,
  commentstyle=\color{cmt}\itshape,
  stringstyle=\color{str},
  literate={\$}{{\textcolor{kw}{\$}}}1,
  aboveskip=8pt, belowskip=8pt,
}
\lstset{style=shell}
\newcommand{\code}[1]{\texttt{\small #1}}

% ---- Headers ----
\pagestyle{fancy}
\fancyhf{}
\lhead{\small Hopper --- Setup \& Architecture}
\rhead{\small\thepage}
\renewcommand{\headrulewidth}{0.4pt}

\titleformat{\section}{\Large\bfseries\color{accent}}{\thesection}{0.6em}{}
\titleformat{\subsection}{\large\bfseries}{\thesubsection}{0.6em}{}

\title{\textbf{Hopper: Local Deployment Journey \\ and Architecture Reference}}
\author{Engineering Notes}
\date{\today}

\begin{document}
\maketitle
\thispagestyle{empty}

\begin{abstract}
\noindent
Hopper is a self-hosted VM-cloud platform for universities: it slices bare-metal
compute into isolated, SSH-accessible containers with credit-based billing and
single sign-on. This document has two goals. First, it is a \emph{field report}:
every concrete problem encountered while bringing the freshly-pulled repository up
on a local machine, the diagnosis, the root cause, and the exact fix. Second, it
is an \emph{architecture reference}: what each moving part (TimescaleDB, Keycloak,
NATS, gRPC, Kubernetes, and the application services) does and \emph{how} it works,
so the next person does not have to re-derive it.
\end{abstract}

\tableofcontents
\newpage

% ===========================================================================
\section{Project Overview}
% ===========================================================================

Hopper lets a university hand every student and professor a lightweight virtual
machine (really a Kubernetes pod) with a fixed CPU/RAM slice, an SSH endpoint, and
a browser-based terminal / VS~Code. Usage is metered in \emph{credits}: an
administrator allocates credits, a professor distributes them to students, and a
running VM burns them every minute until it is stopped or the balance hits zero.

\subsection{High-level architecture}

\begin{lstlisting}
Browser --> Nginx Ingress --> Frontend (SvelteKit)
                          --> API Gateway (FastAPI) --> PostgreSQL (TimescaleDB)
                          --> Keycloak (OIDC SSO)   --> NATS (JetStream)
                                   |
                                   v  gRPC
                             Orchestrator (Go) --> Kubernetes API
                                   |                    |
                                   v                    v
                             Billing Ticker        VM Pods + SSH Services
                                   |
                                   v
                             NATS Events --> API Gateway (credit deduction)
\end{lstlisting}

\subsection{The services at a glance}

\begin{center}
\begin{tabular}{@{}llll@{}}
\toprule
\textbf{Service} & \textbf{Language} & \textbf{Port} & \textbf{Responsibility} \\
\midrule
Frontend     & SvelteKit / Node & 5173 & Dashboard, VM management, credits, admin UI \\
API Gateway  & Python / FastAPI & 8000 & REST API, auth, pods, credits, admin \\
Orchestrator & Go / gRPC        & 50051 & K8s pod lifecycle, billing ticker, metrics \\
PostgreSQL   & TimescaleDB      & 5433 & Users, sessions, double-entry credit ledger \\
NATS         & JetStream        & 4222 & Async events: billing, metrics streaming \\
Keycloak     & Java             & 8080 & OIDC identity provider, university SSO \\
\bottomrule
\end{tabular}
\end{center}

\noindent
Two independent orchestration paths coexist. \textbf{Docker Compose} runs the
stateful infrastructure (PostgreSQL, NATS, Keycloak) locally, while the three
application services run from source. \textbf{Kubernetes} is the production target
and is also where user VMs are scheduled as pods.

% ===========================================================================
\section{Core Technologies Explained}
% ===========================================================================

\subsection{PostgreSQL and TimescaleDB}

\textbf{PostgreSQL} is the relational database of record. It holds users, pod
sessions, and --- most importantly --- the \emph{credit ledger}.

\textbf{TimescaleDB} is a PostgreSQL \emph{extension} (Hopper uses the
\code{timescale/timescaledb} image, which is stock PostgreSQL~16 with the
extension pre-loaded). It adds \emph{hypertables}: a hypertable looks like an
ordinary SQL table but is transparently partitioned into many physical
\emph{chunks} along a time column. This matters because Hopper streams
per-VM CPU/RAM samples: migration \code{008\_metrics\_hypertable} converts the
\code{metrics\_samples} table into a hypertable so that:

\begin{itemize}[nosep]
  \item Inserts of recent time-series data stay fast even as the table grows to
        millions of rows (writes hit the newest chunk only).
  \item Time-range queries (``CPU for pod X over the last hour'') prune to a
        handful of chunks instead of scanning the whole table.
  \item Old data can be dropped or compressed per-chunk cheaply.
\end{itemize}

\noindent
So one engine serves two very different workloads: transactional integrity for
credits, and time-series throughput for metrics.

\subsection{The double-entry credit ledger}

Credits are not a single ``balance'' integer that gets incremented and
decremented --- that is prone to race conditions and leaves no audit trail.
Instead Hopper uses \textbf{double-entry accounting}: every credit movement is a
\emph{transfer} that writes two immutable \code{ledger\_entries} (a debit and a
credit) against two accounts. A user's balance is the \emph{sum} of their
entries. A special \emph{system account} (\code{00000000-\dots-0000}) is the
counterparty for allocations and deductions. PostgreSQL \emph{advisory locks}
serialise concurrent deductions on the same account so a VM ticking every minute
can never double-spend or race an admin allocation.

\subsection{Keycloak (identity, SSO, OIDC)}

\textbf{Keycloak} is an open-source identity and access-management server. It is
the single source of truth for \emph{who a user is} and \emph{what role they
have}. Hopper never stores passwords itself.

Key concepts:
\begin{itemize}[nosep]
  \item \textbf{Realm} --- an isolated tenant containing its own users, roles,
        and clients. Hopper uses a realm named \code{hopper}.
  \item \textbf{Client} --- an application that asks Keycloak to authenticate
        users. Hopper has two:
        \begin{itemize}[nosep]
          \item \code{hopper-api} --- a \emph{public} client used by the browser
                login (authorization-code + PKCE, and direct password grant for
                the in-app themed login).
          \item \code{hopper-admin} --- a \emph{confidential} service-account
                client the API gateway uses (via the \emph{client-credentials}
                grant) to mutate Keycloak: create users on signup, mark emails
                verified, reset passwords, change roles.
        \end{itemize}
  \item \textbf{Roles} --- \code{admin}, \code{professor}, \code{student}. They
        travel inside the JWT under \code{realm\_access.roles}.
  \item \textbf{OIDC / OAuth 2.1} --- the login uses \textbf{PKCE} (proof key for
        code exchange) so an intercepted authorization code is useless without
        the verifier, plus \code{state} (CSRF) and \code{nonce} (replay)
        protection.
\end{itemize}

\noindent
\textbf{How a login works.} The browser hits \code{GET /auth/login}; the API
gateway redirects it to Keycloak's hosted page with a PKCE challenge. After the
user authenticates, Keycloak redirects back to
\code{/api/auth/callback} with a one-time code; the gateway exchanges the code
(plus the PKCE verifier) for a signed JWT, validates the token's signature
(against Keycloak's JWKS), issuer, expiry, and \emph{audience}, then sets an
\code{HttpOnly} session cookie. The token's \code{aud} claim \emph{must} include
\code{hopper-api} or the gateway rejects it --- this detail caused a real bug,
covered in Section~\ref{sec:problems}.

\subsection{NATS and JetStream (the message bus)}

\textbf{NATS} is a lightweight, high-throughput \emph{publish/subscribe} message
broker. Instead of services calling each other synchronously for every event,
they publish messages to named \emph{subjects}; interested services subscribe.
This decouples producers from consumers.

\textbf{JetStream} is NATS's persistence layer. Plain NATS is fire-and-forget (a
message with no listener is lost); JetStream adds durable streams so a message is
stored and redelivered until acknowledged --- essential for billing, where a lost
``deduct credits'' event would mean free compute.

\noindent
\textbf{How Hopper uses it:}

\begin{center}
\begin{tabular}{@{}llp{6.3cm}@{}}
\toprule
\textbf{Subject} & \textbf{Publisher $\rightarrow$ Consumer} & \textbf{Purpose} \\
\midrule
\code{pod.created}       & Orchestrator $\rightarrow$ --- & Pod lifecycle event \\
\code{pod.stopped}       & Orchestrator $\rightarrow$ --- & Pod terminated \\
\code{billing.deducted}  & Orchestrator $\rightarrow$ API GW & Deduct credits from the ledger \\
\code{billing.exhausted} & API GW $\rightarrow$ Orchestrator & Auto-terminate a VM at zero balance \\
\code{metrics.<pod\_id>} & Orchestrator $\rightarrow$ API GW & Real-time CPU/RAM samples (SSE) \\
\bottomrule
\end{tabular}
\end{center}

\noindent
\textbf{The billing loop}, concretely: the orchestrator's ticker fires every 60s
for each running VM and publishes \code{billing.deducted}. The API gateway
consumes it and calls \code{deduct\_credits()} on the ledger. If the balance
reaches zero, the gateway publishes \code{billing.exhausted}; the orchestrator
consumes that and terminates the pod. Metrics flow the same way: the orchestrator
samples every 5s, publishes to \code{metrics.<pod\_id>}, and the gateway relays
them to the browser as \textbf{Server-Sent Events (SSE)}.

\subsection{gRPC and Protocol Buffers}

The API gateway (Python) and the orchestrator (Go) talk over \textbf{gRPC}, a
binary RPC protocol. The contract is defined once in \code{.proto} files
(\code{proto/hopper/pod/v1/pod.proto},
\code{proto/hopper/billing/v1/billing.proto}) and \emph{code-generated} into both
languages: Python stubs for the gateway, Go stubs for the orchestrator. The
gateway calls typed methods --- \code{CreatePod}, \code{TerminatePod},
\code{GetPodStatus}, \code{StreamMetrics}, \code{ListNodes} --- as if they were
local functions. Because the stubs are generated, editing a \code{.proto} without
regenerating leaves the two sides out of sync; this also caused a real build
failure (Section~\ref{sec:problems}).

\subsection{Kubernetes (the compute substrate)}

A user ``VM'' is a Kubernetes \textbf{pod} with CPU/memory \emph{limits} matching
the chosen plan (small/medium/large), fronted by a \textbf{NodePort Service}
(ports 30000--32767) that exposes SSH. The orchestrator is the only component
that talks to the Kubernetes API: it creates the pod, watches its lifecycle,
publishes metrics, and tears it down. In production, \textbf{Cilium} network
policies enforce zero-trust isolation so a user VM can reach the internet but not
other VMs or the platform's internal services.

\subsection{The application services}

\begin{itemize}[nosep]
  \item \textbf{Frontend (SvelteKit)} --- the web UI, served in dev by Vite on
        port 5173. It proxies \code{/api/*} to the gateway.
  \item \textbf{API Gateway (FastAPI)} --- the REST surface and security boundary:
        auth, pod CRUD, credits, admin endpoints, SSE metric relay, and the NATS
        billing/metrics consumers. Migrations are managed with \textbf{Alembic}.
  \item \textbf{Orchestrator (Go)} --- the gRPC server, the Kubernetes client,
        the billing ticker, and the metrics publisher.
\end{itemize}

% ===========================================================================
\section{Running the Project: Problems and Solutions}
\label{sec:problems}
% ===========================================================================

The repository had been updated since a prior successful run. The update moved
the PostgreSQL port, added nine database migrations, introduced new frontend and
backend dependencies, changed the proto contract, and (implicitly) required a
provisioned Keycloak realm. Bringing it up surfaced \textbf{seven} distinct
blockers. Each is documented below with symptom, diagnosis, root cause, and fix.

\subsection*{Environment}
Linux, Docker Compose for infra, host had: Docker, a system PostgreSQL install,
Node~24 (via nvm) with pnpm~9 system-wide, Go, and the protobuf toolchain.

% ------------------------------------------------------------------
\subsection{Problem 1 --- PostgreSQL port 5433 already in use}

\textbf{Symptom.} \code{docker compose up} failed:
\begin{lstlisting}
Error response from daemon: ... failed to bind host port 0.0.0.0:5433/tcp:
address already in use
\end{lstlisting}

\textbf{Diagnosis.} The updated \code{docker-compose.yml} publishes PostgreSQL on
host port \textbf{5433} (the previous version used 5434). Something already owned
5433:
\begin{lstlisting}
$ ss -ltn | grep :5433
LISTEN 0 200 127.0.0.1:5433 0.0.0.0:*
$ pg_lsclusters
16  main  5433 online postgres ...
\end{lstlisting}

\textbf{Root cause.} A \emph{host-installed} PostgreSQL~16 cluster
(\code{postgresql@16-main}, a systemd service) was squatting on exactly the port
the project now wants for its Docker PostgreSQL. The application config
(\code{app/config.py}) and \code{alembic.ini} both hard-code \code{localhost:5433},
so remapping the container alone would not have been enough.

\textbf{Fix.} Stop (and disable, so it does not return on reboot) the host
cluster, freeing the port for Docker:
\begin{lstlisting}
$ sudo systemctl stop postgresql@16-main
$ sudo systemctl disable postgresql@16-main   # so a reboot doesn't reclaim 5433
\end{lstlisting}
It recurred twice because a plain \code{stop} does not survive a reboot ---
hence the \code{disable}.

% ------------------------------------------------------------------
\subsection{Problem 2 --- pnpm version mismatch}

\textbf{Symptom.} The frontend dependency install aborted:
\begin{lstlisting}
ERROR  packages field missing or empty
\end{lstlisting}

\textbf{Diagnosis \& root cause.} The update added
\code{frontend/pnpm-workspace.yaml} containing only an \code{allowBuilds:} key ---
a \textbf{pnpm~10+} feature that whitelists native build scripts (esbuild's, in
this case). Under pnpm~10+, a \code{pnpm-workspace.yaml} may omit the
\code{packages:} field; under \textbf{pnpm~9} (what the host had, system-wide at
\code{/usr/local/bin}) that field is mandatory, so pnpm~9 rejects the file
outright. The repo now assumes pnpm~10/11 locally.

\textbf{Fix.} Install pnpm~11 without touching the root-owned system binary or
requiring \code{sudo}, using Corepack (bundled with Node) into a user-writable
directory, then put it first on \code{PATH}:
\begin{lstlisting}
$ corepack prepare pnpm@latest --activate
$ corepack enable --install-directory "$HOME/.local/bin" pnpm
$ export PATH="$HOME/.local/bin:$PATH"
$ pnpm --version      # -> 11.x  (system pnpm 9 left untouched)
\end{lstlisting}

% ------------------------------------------------------------------
\subsection{Problem 3 --- pnpm aborts without a TTY}

\textbf{Symptom.}
\begin{lstlisting}
ERR_PNPM_ABORTED_REMOVE_MODULES_DIR_NO_TTY
Aborted removal of modules directory due to no TTY
\end{lstlisting}

\textbf{Root cause.} A \code{node\_modules} left behind by pnpm~9 was
incompatible with pnpm~11, which wanted to purge and rebuild it --- but purging
asks for interactive confirmation, and the launcher ran headless (background,
no TTY).

\textbf{Fix.} Signal non-interactive/CI mode and pre-clean the stale directory:
\begin{lstlisting}
$ rm -rf frontend/node_modules
$ export CI=true          # pnpm auto-confirms the purge
\end{lstlisting}

% ------------------------------------------------------------------
\subsection{Problem 4 --- the whole stack tore itself down between steps}

\textbf{Symptom.} Repeatedly, the background run would end (exit 137/143) and
\emph{every} container and service would be gone.

\textbf{Root cause.} \code{scripts/deploy-local.sh} runs a foreground supervisor
loop and installs a \code{trap} on \code{SIGTERM/SIGINT} whose cleanup runs
\code{docker compose down} and kills the app children. Whenever the controlling
background task was stopped, that trap fired and demolished the stack.

\textbf{Fix.} Detach the supervisor into its own session with \code{setsid} so a
stop of the controlling task does not deliver the signal that triggers teardown:
\begin{lstlisting}
$ setsid ./scripts/deploy-local.sh > .local-logs/deploy.out 2>&1 &
\end{lstlisting}
A related trap: \code{pkill -f 'deploy-local.sh'} matches the \emph{pkill command
line itself} and self-terminates the shell (exit 144); stop the supervisor by PID
instead.

% ------------------------------------------------------------------
\subsection{Problem 5 --- orchestrator would not compile (stale proto stubs)}

\textbf{Symptom.} The Go service exited immediately:
\begin{lstlisting}
internal/grpc/pod_service.go:57: unknown field SshPassword in struct literal
   of type podv1.PodStatus
internal/grpc/pod_service.go:100: req.PodId undefined
   (type *podv1.CreatePodRequest has no field or method PodId)
\end{lstlisting}

\textbf{Diagnosis.} The \code{.proto} files had gained \code{ssh\_password} and
\code{pod\_id} fields, but the committed \emph{generated} Go stubs were stale ---
they predated those fields. A downstream effect: the API gateway hung on
``connecting to orchestrator\dots'' (and never opened port 8000) because the
orchestrator it depends on never came up.

\textbf{Root cause \& a second bug.} The regeneration script
(\code{scripts/generate-proto.sh}) prefers \code{buf} if installed, but
\code{proto/buf.gen.yaml} was \textbf{empty}, which breaks the \code{buf} path.
The \code{protoc} fallback, however, was fully functional. Note also the layout
subtlety: the code imports
\code{.../api/proto/hopper/pod/v1}, and only \code{protoc} with
\code{paths=source\_relative} reproduces that exact directory.

\textbf{Fix.} Regenerate the Go stubs with \code{protoc} directly (Python stubs
were already current):
\begin{lstlisting}
$ export PATH="$HOME/go/bin:$PATH"   # protoc-gen-go, protoc-gen-go-grpc
$ GO_OUT=services/orchestrator/api/proto
$ rm -rf "$GO_OUT" && mkdir -p "$GO_OUT"
$ protoc -Iproto \
    --go_out="$GO_OUT" --go_opt=paths=source_relative \
    --go-grpc_out="$GO_OUT" --go-grpc_opt=paths=source_relative \
    proto/hopper/pod/v1/pod.proto proto/hopper/billing/v1/billing.proto
$ (cd services/orchestrator && go build ./cmd/orchestrator/)   # -> BUILD OK
\end{lstlisting}

% ------------------------------------------------------------------
\subsection{Problem 6 --- database migrations}

Not an error so much as a required step: the update added \textbf{nine} Alembic
migrations (\code{004}--\code{012}: VS~Code port, audit logs, ssh keys, billing
accuracy, the metrics hypertable, ssh password, issue reports, the
pending-teacher flag, and email verification codes). The deploy script applies
them automatically (\code{alembic upgrade head}); a manual run is:
\begin{lstlisting}
$ cd services/api-gateway && PYTHONPATH=. poetry run alembic upgrade head
\end{lstlisting}

% ------------------------------------------------------------------
\subsection{Problem 7 --- Keycloak: ``We are sorry\dots unable to login''}

\textbf{Symptom.} The Keycloak-hosted sign-in page showed a generic
``We are sorry\dots'' error and login was impossible.

\textbf{Diagnosis.} The app redirects the browser to
\code{.../realms/hopper/protocol/openid-connect/auth}, but that realm returned
404:
\begin{lstlisting}
$ curl -s -o /dev/null -w '%{http_code}' \
    http://localhost:8080/realms/hopper/.well-known/openid-configuration
404
\end{lstlisting}

\textbf{Root cause.} \code{deploy-local.sh} starts Keycloak but never
\emph{provisions} it. A fresh Keycloak has only the \code{master} realm, so the
\code{hopper} realm, its clients, roles, and users simply did not exist. (The
repo's only Keycloak script, \code{keycloak-harden.sh}, targets the Kubernetes
deployment via \code{kubectl exec}, not the local container.)

\textbf{Fix.} Provision the realm \emph{idempotently} with \code{kcadm} inside the
container. The essential objects:

\begin{itemize}[nosep]
  \item Realm \code{hopper} --- \code{sslRequired=NONE} (local HTTP),
        registration enabled.
  \item Realm roles \code{admin}, \code{professor}, \code{student}.
  \item Public client \code{hopper-api} --- standard + direct-access grants,
        redirect URI \code{http://localhost:5173/api/auth/callback}, web origins,
        PKCE, and \textbf{an audience mapper adding \code{hopper-api} to the token
        \code{aud}}. This last item is easy to miss and \emph{critical}:
        \code{middleware/auth.py} validates \code{audience=hopper-api}, and
        Keycloak does not add the client to its own \code{aud} by default ---
        without the mapper, \emph{every} login succeeds at Keycloak but the
        gateway rejects the resulting token.
  \item Confidential client \code{hopper-admin} --- service account with
        \code{manage-users}/\code{view-users} and a shared secret, for
        signup / verification / reset / role changes.
  \item A ready-to-use user \code{admin@hopper.dev} (email verified, \code{admin}
        role).
\end{itemize}

\noindent
The admin client's secret is handed to the gateway via a gitignored
\code{services/api-gateway/.env}:
\begin{lstlisting}
HOPPER_KEYCLOAK_ADMIN_CLIENT_SECRET=hopper-admin-local-secret
\end{lstlisting}

\textbf{Verification.} After provisioning, a password-grant login returned a token
with exactly the right claims:
\begin{lstlisting}
aud        = ['hopper-api', 'account']   # includes hopper-api  -> gateway accepts
azp        = hopper-api
iss        = http://localhost:8080/realms/hopper
email      = admin@hopper.dev | verified = True
realm_roles= ['admin']
\end{lstlisting}
and through the gateway: \code{POST /auth/login} $\rightarrow$ \textbf{200}
(\code{role: admin}); \code{GET /auth/login} $\rightarrow$ \textbf{307} (redirect
to Keycloak). Because Keycloak is Postgres-backed, the realm survives restarts and
\code{docker compose down}; only a \code{-v} volume wipe loses it.

% ===========================================================================
\section{Final Working Setup and Verification}
% ===========================================================================

With all seven blockers resolved, the full launch (dependencies already installed
and migrations already applied) is:
\begin{lstlisting}
$ export COREPACK_ENABLE_DOWNLOAD_PROMPT=0 CI=true
$ export PATH="$HOME/.local/bin:$HOME/.nvm/versions/node/v24.12.0/bin:\
/snap/go/current/bin:$HOME/go/bin:$PATH"
$ setsid ./scripts/deploy-local.sh --skip-install --skip-migrate \
    > .local-logs/deploy.out 2>&1 &
\end{lstlisting}

\noindent
Health checks, all green:

\begin{center}
\begin{tabular}{@{}lll@{}}
\toprule
\textbf{Check} & \textbf{Expected} & \textbf{Result} \\
\midrule
\code{docker compose ps}            & postgres/nats/keycloak Up & healthy \\
API \code{/docs}                    & 200 & 200 \\
Frontend \code{/}                   & 200/302 & 302 (login) \\
Keycloak \code{/realms/hopper}      & 200 & 200 \\
\code{POST /auth/login}             & 200 & 200, role=admin \\
Orchestrator \code{:50051}          & listening & listening \\
Alembic head                        & 012 & 012 \\
\bottomrule
\end{tabular}
\end{center}

% ===========================================================================
\section{Quick Reference}
% ===========================================================================

\subsection{Endpoints and credentials (local)}

\begin{center}
\begin{tabular}{@{}lll@{}}
\toprule
\textbf{Service} & \textbf{URL} & \textbf{Credentials} \\
\midrule
Frontend      & \code{http://localhost:5173}       & \code{admin@hopper.dev} / \code{Admin123!} \\
API docs      & \code{http://localhost:8000/docs}  & --- \\
Keycloak      & \code{http://localhost:8080}       & \code{admin} / \code{admin} \\
PostgreSQL    & \code{localhost:5433}              & \code{hopper} / \code{hopper\_dev} \\
NATS monitor  & \code{http://localhost:8222}       & --- \\
\bottomrule
\end{tabular}
\end{center}

\subsection{Lessons distilled}

\begin{enumerate}[nosep]
  \item A moved host port (5433) can be squatted by a \emph{system} service ---
        check \code{ss}/\code{pg\_lsclusters}, not just Docker.
  \item Tool-version drift is real: a \code{pnpm-workspace.yaml} written for
        pnpm~10+ hard-fails on pnpm~9. Corepack pins versions without sudo.
  \item Headless runs need \code{CI=true}; long-lived supervisors need
        \code{setsid} to survive; never \code{pkill -f} a pattern that matches
        your own command.
  \item Generated code (proto stubs) must be regenerated when the contract
        changes; \code{paths=source\_relative} must match the import layout.
  \item Keycloak is not usable until a realm is provisioned --- and the
        token \code{aud} must include the API's client ID or every token is
        rejected downstream.
\end{enumerate}

\end{document}
